\chapter{Does It Transfer? --- The Four-Cell Question}
\label{ch:transfer}

\begin{plainlanguage}
\textbf{In plain terms:} If recognition and implementation really are separate skills, does practicing one help the other? Neither flat independence nor full unity survives the evidence. The two are clearly distinguishable but not cleanly separate — entangled, not identical.
\end{plainlanguage}

\section{Why transfer is a stronger test than performance}

Before the matrix machinery, it is worth being explicit about why this chapter asks about transfer at all rather than stopping at performance, since the move from one to the other is not obvious to a reader encountering it for the first time. A competence that only succeeds in the exact environment it was trained or observed in is weaker evidence of a real, general capacity than one that succeeds somewhere new --- because the first case leaves open a deflationary explanation the second case rules out. A child who reliably refuses the basket-for-water counterfoil might be pattern-matching a specific, memorized scenario rather than possessing any general competence in admissibility recognition; a child who recognizes an unfamiliar, never-before-seen counterfoil using the same underlying logic is harder to explain away this way. Performance in the training environment is therefore necessary but not sufficient evidence for the kind of competence this book is interested in tracking. Transfer --- performance somewhere the training did not directly prepare for --- is the test that actually discriminates a genuine, generalizable capacity (the \emph{competence} row of the previous chapter's table) from the case where correct behavior was produced some other way (the \emph{luck or habit} row). This is why the four-cell question below matters more than it might first appear: it is not an incidental refinement of the recognition/implementation split, it is the test that determines whether the split was tracking something real in the first place.

\section{The transfer matrix}

If recognition and implementation are genuinely distinct competencies, the natural next question is whether training in one transfers to performance in the other, and the question splits into four parts rather than one.

\begin{definition}[Transfer Matrix]
\[
T = \begin{pmatrix} T_{RR} & T_{RI} \\ T_{IR} & T_{II} \end{pmatrix}
\]
where $T_{RR}$ measures whether recognition practice at high margins improves recognition at lower margins; $T_{II}$ measures whether implementation practice under one temptation improves implementation under another; $T_{IR}$ measures whether implementation practice improves later recognition, particularly recognition under pressure; and $T_{RI}$ measures whether recognition practice improves later implementation.
\end{definition}

\begin{naivemodel}
A first pass through this question risks an overly tidy answer: strong within-faculty transfer ($T_{RR}, T_{II} > 0$) alongside near-zero cross-faculty transfer ($T_{RI} \approx T_{IR} \approx 0$), which would license a clean story of two largely independent developmental pathways --- the familiar folk distinction between being ``brilliant but self-destructive'' versus ``wise but ineffective,'' reframed as two points in an $(I_R, I_I)$ plane rather than as a contradiction.
\end{naivemodel}

\begin{proposition}[Orthogonality Failure]
If $T_{IR} > 0$ or $T_{RI} > 0$, then $I_R$ and $I_I$ are not orthogonal in the relevant sense.
\end{proposition}

This is true essentially by the definition of what nonzero cross-transfer means, and is stated for clarity rather than as a result requiring independent proof. Whether the antecedent holds --- whether $T_{IR}$ or $T_{RI}$ is actually nonzero --- is the empirical question this chapter exists to address.

\section{What the literature shows}

The clean independence story does not survive contact with the developmental and social-psychology literature, and it fails in an informative rather than merely disappointing way. The construct of \emph{moral disengagement} --- the use of self-monitoring, judgment, and self-reactive mechanisms to permit wrongdoing one already recognizes as wrong --- maps closely onto the high-$I_R$, low-$I_I$ cell, and its existence as an independently studied construct is real evidence that recognition and implementation dissociate. But a longitudinal finding complicates flat independence in the other direction: self-control and cooperative behavior in childhood predict reduced moral disengagement in adolescence --- early implementation-relevant practice associated with a later, partly recognition-adjacent outcome. This is correlational rather than experimental evidence, and the direction of any underlying causal mechanism should be stated as a hypothesis the finding is consistent with, not a demonstrated mechanism --- but it is a documented nonzero association, inconsistent with a flat $T_{IR} \approx 0$ claim. Separately, the broader self-regulation literature warns against treating $I_I$ as a single transferable competence at all: executive function, emotion regulation, response inhibition, and effortful control are measurably distinct, overlapping but non-identical constructs, which means even $T_{II}$ --- transfer within implementation alone --- should not be assumed to hold across arbitrarily different forms of temptation.

\begin{proposition}[Dissociability]
If there exist population subsets in which $I_R \gg I_I$, and other subsets in which $I_I \gg I_R$, then the two capacities are non-identical.
\end{proposition}

Also near-definitional --- the existence of such subsets is itself the content of the claim --- and supported here by citing literature consistent with both subset types existing (moral disengagement for the first; documented cases of strong impulse control without corresponding insight for the second) rather than derived mathematically.

\section{The honest formal conclusion}

\begin{repairbox}
\[
I_R \not\equiv I_I \qquad \text{but} \qquad I_R \not\perp I_I
\]

Recognition and implementation are dissociable as constructs, not identical --- but they are not cleanly independent either. They are partially coupled, developmentally entangled in at least one demonstrated direction, and --- as the self-regulation literature insists --- individually multidimensional rather than scalar. This last point reframes $I_I$ itself:
\[
I_I = (I_{\text{inhibition}}, I_{\text{emotion}}, I_{\text{effort}}, I_{\text{delay}}, I_{\text{social-cost}}, \ldots)
\]
a vector of partially correlated sub-competencies rather than a single number. Implementation, offered in Chapter~\ref{ch:two-axes} as the second half of a clean two-term split, turns out under inspection to be its own compressed bundle. This is the observation Chapter~\ref{ch:recursion} exists to take seriously rather than treat as an inconvenience.
\end{repairbox}
